\documentclass[aspectratio=169]{beamer}

\usepackage{graphicx}
\usepackage{amsmath,amssymb}
\usepackage{ragged2e}
\usepackage{xcolor}
\usepackage{tikz}
\usepackage{array}

\graphicspath{{./}{Figures/}{MACC2026/Figures/}}

\definecolor{maroon1}{rgb}{0.478, 0, 0.235}
\definecolor{lightgray}{rgb}{0.94,0.94,0.97}
\definecolor{dimgray}{rgb}{0.369, 0.416, 0.443}
\definecolor{maccblue}{RGB}{28,74,117}
\definecolor{maccgreen}{RGB}{44,118,92}
\definecolor{maccred}{RGB}{156,58,52}
\definecolor{softblue}{RGB}{236,243,250}
\definecolor{softgreen}{RGB}{235,247,242}
\definecolor{softred}{RGB}{251,239,238}



\usetheme{Copenhagen}
\usecolortheme[named=maroon1]{structure}

\setbeamerfont{caption}{size=\scriptsize}
\setbeamertemplate{navigation symbols}{}
\setbeamercolor{footerleft}{fg=white,bg=black}
\setbeamercolor{footerright}{fg=white,bg=maroon1}
\setbeamertemplate{footline}{%
  \leavevmode%
  \hbox{%
    \begin{beamercolorbox}[wd=0.5\paperwidth,ht=1.95ex,dp=0.65ex,leftskip=0.35cm,rightskip=0.15cm]{footerleft}%
      \color{white}\tiny Amir Hamedi
    \end{beamercolorbox}%
    \begin{beamercolorbox}[wd=0.5\paperwidth,ht=1.95ex,dp=0.65ex,leftskip=0.15cm,rightskip=0.35cm]{footerright}%
      \color{white}\tiny \hfill MACC Annual Meeting 2026
    \end{beamercolorbox}%
  }%
}
\setbeamercolor{block body}{fg=black,bg=lightgray}
\setbeamercolor{block title}{fg=white,bg=maroon1}

\title[MACC 2026]{A Practical Reinforcement Learning Control Framework for Complex Process Systems}
\author[A. H. Hamedi]{A. H. Hamedi}
\date{}

\newcommand{\PosterTitle}[1]{%
  \frametitle{\centering \scriptsize \textbf{#1}}%
  \vspace{-0.48cm}%
}

\newcommand{\ColumnHeader}[2]{%
  {\setlength{\fboxsep}{2pt}\setlength{\fboxrule}{0pt}%
  \noindent\colorbox{#1}{%
    \begin{minipage}[c][0.047\textheight][c]{\dimexpr\linewidth-2\fboxsep\relax}
      \centering\color{white}\bfseries\scriptsize #2
    \end{minipage}}}%
}

\newcommand{\FigurePanel}[2]{%
  {\setlength{\fboxsep}{1pt}\setlength{\fboxrule}{0.35pt}%
  \noindent\fcolorbox{black!25}{white}{%
    \begin{minipage}[c][#2][c]{\dimexpr\linewidth-2\fboxsep-2\fboxrule\relax}
      \centering
      \includegraphics[width=0.98\linewidth,height=#2,keepaspectratio]{#1}
    \end{minipage}}}%
}

\newcommand{\TextPanel}[4]{%
  {\setlength{\fboxsep}{2pt}\setlength{\fboxrule}{0.55pt}%
  \noindent\fcolorbox{#1}{#2}{%
    \begin{minipage}[t][#3][t]{\dimexpr\linewidth-2\fboxsep-2\fboxrule\relax}
      #4
    \end{minipage}}}%
}

\newcommand{\ResultPanel}[3]{%
  {\setlength{\fboxsep}{2pt}\setlength{\fboxrule}{0.55pt}%
  \noindent\fcolorbox{#1}{#2}{%
    \begin{minipage}[c][0.17\textheight][c]{\dimexpr\linewidth-2\fboxsep-2\fboxrule\relax}
      \centering #3
    \end{minipage}}}%
}

\newcommand{\WideResultPanel}[3]{%
  {\setlength{\fboxsep}{2pt}\setlength{\fboxrule}{0.55pt}%
  \noindent\fcolorbox{#1}{#2}{%
    \begin{minipage}[c][0.17\textheight][c]{\dimexpr\linewidth-2\fboxsep-2\fboxrule\relax}
      \centering #3
    \end{minipage}}}%
}

\newcommand{\SmallBullet}[1]{%
  \item \justifying #1
}

\begin{document}

%=========================================================
\begin{frame}
\frametitle{%
  \hspace{0.02cm}\scriptsize\textbf{Practical Reinforcement Learning for Process Control}%
  \hfill
  \scriptsize\textbf{Amir Hamedi \hspace{0.6em}$|$\hspace{0.6em} Dr.~Mhaskar}%
}
\vspace{-0.28cm}
\scriptsize

\vspace{0.08cm}

\begin{block}{\textbf{Why RL for Industrial Process Control?}}
\begin{itemize}\justifying
  \item Industrial MPC is often built from linear identified models; performance can degrade when nonlinear dynamics, coupling, or operating conditions change.
  \item RL can adapt online from closed-loop data, but unrestricted exploration is not acceptable for industrial plants.
\end{itemize}
\end{block}



\begin{columns}[T,onlytextwidth]
\begin{column}{0.32\textwidth}
\ColumnHeader{maccblue}{Why MPC alone is not enough?}
\vspace{0.06cm}
\TextPanel{maccblue}{softblue}{0.53\textheight}{%
\begin{itemize}
  \SmallBullet{Linear MPC is practical, interpretable, and widely used.}
  \SmallBullet{A fixed linear model can lose accuracy across nonlinear operating regions.}
  \SmallBullet{Reidentification and retuning are expensive when plant behavior drifts.}
  \SmallBullet{This motivates controllers that adapt without discarding MPC.}
\end{itemize}
}
\end{column}

\begin{column}{0.32\textwidth}
\ColumnHeader{maccgreen}{RL learns from interaction}
\vspace{0.06cm}
\TextPanel{maccgreen}{softgreen}{0.53\textheight}{%
\centering
\includegraphics[width=0.98\linewidth,height=0.2\textheight,keepaspectratio]{RL Flowdiagram.png}

\vspace{0.05cm}
\begin{itemize}
  \SmallBullet{RL improves through repeated plant interaction.}
  \SmallBullet{Naive online training requires extensive exploration.}
  \SmallBullet{This can cause unsafe moves and product-quality excursions.}

\end{itemize}
}
\end{column}

\begin{column}{0.32\textwidth}
\ColumnHeader{maccred}{Practical RL control approaches}
\vspace{0.06cm}
\TextPanel{maccred}{softred}{0.53\textheight}{%
\begin{itemize}
  \SmallBullet{\textbf{Pretraining:} initialize RL from MPC behavior before deployment.}
  \SmallBullet{\textbf{RL-assisted MPC:} let RL tune MPC variables instead of directly controlling the plant.}
  \SmallBullet{\textbf{Safe RL:} filter or replace unsafe RL actions using model-based checks.}
\end{itemize}

\vspace{0.16cm}
\centering
{\color{maccred}\bfseries Keep MPC in the loop}\\[-0.03em]
}
\end{column}
\end{columns}

\end{frame}

%=========================================================
\begin{frame}
\PosterTitle{Pretrained RL (Published in I\&EC Research)}
\scriptsize

\vspace{0.25cm}


\begin{columns}[T,onlytextwidth]
\begin{column}{0.57\textwidth}
\ColumnHeader{maccblue}{MPC pretraining and online refinement}
\vspace{0.06cm}
\FigurePanel{PretrainRL.png}{0.62\textheight}
\end{column}

\begin{column}{0.40\textwidth}
\ColumnHeader{maccgreen}{Previous approach limitation}
\vspace{0.06cm}
\TextPanel{maccgreen}{softgreen}{0.25\textheight}{%
\begin{itemize}
  \SmallBullet{Previous work used the MPC one-step quadratic reward during online training.}
  \SmallBullet{That reward reduced unsafe exploration, but still left significant steady-state error.}
\end{itemize}
}

\vspace{0.055cm}
\ColumnHeader{maccblue}{Modifications}
\vspace{0.06cm}
\TextPanel{maccblue}{softblue}{0.25\textheight}{%
\begin{itemize}
  \SmallBullet{Offline TD3 initialization imitates offset-free MPC trajectories.}
  \SmallBullet{Online reward adds offset-aware shaping near the setpoint.}
  \SmallBullet{Mixed replay balances priority, recent data, and broad coverage.}
\end{itemize}
}
\end{column}
\end{columns}

\vspace{-0.08cm}
\WideResultPanel{maccblue}{softblue}{%
{\color{maccblue}\bfseries Our approach = MPC imitation + offset-aware online TD3 + mixed replay}\\[-0.04em]
{Same practical starting point as MPC-pretrained RL, but with stronger learning signal for near offset-free behavior.}
}
\end{frame}

%=========================================================
\begin{frame}
\PosterTitle{Pretrained RL Results: Polymer Reactor and Aspen Dynamics C$_2$ Splitter}
\scriptsize

\vspace{0.2cm}

\begin{columns}[T,onlytextwidth]
\begin{column}{0.59\textwidth}
\ColumnHeader{maccblue}{Industrial C$_2$ implementation}
\vspace{0.04cm}
\FigurePanel{c2_interface_td3.png}{0.35\textheight}

% \vspace{0.025cm}
% \TextPanel{maccblue}{softblue}{0.025\textheight}{%
% \centering
% {\color{maccblue}\bfseries \tiny Aspen Dynamics + Python TD3}\\[-0.08em]
% }

\end{column}

\begin{column}{0.39\textwidth}
\ColumnHeader{maccgreen}{Polymer reactor: final offset reduction}
\vspace{0.06cm}
\TextPanel{maccgreen}{softgreen}{0.12\textheight}{%
\centering
\tiny
\setlength{\tabcolsep}{3.4pt}
\renewcommand{\arraystretch}{0.96}
\begin{tabular}{lcccc}
\hline
 & \multicolumn{2}{c}{Nominal} & \multicolumn{2}{c}{Oper. changes} \\
Output & SP(a) & SP(b) & SP(a) & SP(b) \\
\hline
\(\eta\) & \textbf{98\%} & \textbf{90\%} & \textbf{92\%} & \textbf{76\%} \\
\(T\) & \textbf{92\%} & \textbf{84\%} & \textbf{84\%} & \textbf{90\%} \\
\hline
\end{tabular}\\[0.2em]
}

% \vspace{0.05cm}
\ColumnHeader{maccred}{C$_2$ splitter: final offset reduction}
\vspace{0.06cm}
\TextPanel{maccred}{softred}{0.12\textheight}{%
\centering
\tiny
\setlength{\tabcolsep}{2.15pt}
\renewcommand{\arraystretch}{0.92}
\begin{tabular}{lcccccc}
\hline
 & \multicolumn{2}{c}{Nom.} & \multicolumn{2}{c}{Fluct.} & \multicolumn{2}{c}{Ramp} \\
Output & SP1 & SP2 & SP1 & SP2 & SP1 & SP2 \\
\hline
\(x_{24}\) & \textbf{99.2} & \textbf{99.6} & \textbf{99.2} & \textbf{89.7} & \textbf{94.9} & \textbf{91.1} \\
\(T_{85}\) & \textbf{81.3} & \textbf{90.9} & \textbf{91.6} & \textbf{60.9} & \textbf{35.8} & \textbf{97.8} \\
\hline
\end{tabular}\\[0.2em]
}
\end{column}
\end{columns}

\vspace{0.035cm}
\ColumnHeader{maccblue}{Last evaluation episodes: C$_2$ splitter comparisons}
\vspace{-0.25cm}
\begin{columns}[T,onlytextwidth]
\begin{column}{0.32\textwidth}
\centering
\FigurePanel{C2_nom_outputs_compare_three.png}{0.35\textheight}\\[-0.1em]
{\tiny Nominal}
\end{column}
\begin{column}{0.32\textwidth}
\centering
\FigurePanel{C2_rampnoise_outputs_compare_three.png}{0.35\textheight}\\[-0.1em]
{\tiny Fluctuating operation}
\end{column}
\begin{column}{0.32\textwidth}
\centering
\FigurePanel{C2_ramp_outputs_compare_three.png}{0.35\textheight}\\[-0.1em]
{\tiny Ramp operation}
\end{column}
\end{columns}
\end{frame}

%=========================================================
\begin{frame}
\PosterTitle{RL-Assisted MPC}
\scriptsize

\begin{block}{\textbf{Main idea}}
\begin{itemize}\justifying
  \item RL assists the existing controller instead of replacing it.
  \item MPC still solves the constrained control problem; RL tunes selected MPC design variables.
\end{itemize}
\end{block}

\vspace{-0.16cm}
\begin{columns}[T,onlytextwidth]
\begin{column}{0.58\textwidth}
\ColumnHeader{maccgreen}{RL-assisted MPC architecture}
\vspace{0.06cm}
\FigurePanel{RLAssistedDiagram.png}{0.52\textheight}
\end{column}

\begin{column}{0.39\textwidth}
\ColumnHeader{maccgreen}{Case studies}
\vspace{0.045cm}
\TextPanel{maccgreen}{softgreen}{0.1\textheight}{%
\begin{itemize}
  \SmallBullet{Polymer CSTR.}
  \SmallBullet{Aspen Dynamics C$_2$ splitter.}
\end{itemize}
}

\vspace{0.045cm}
\ColumnHeader{maccblue}{What the agents change in MPC}
\vspace{0.045cm}
\TextPanel{maccblue}{softblue}{0.35\textheight}{%
\begin{itemize}
  \SmallBullet{\textbf{Penalty weights:} adapt tracking and input-move}
  \SmallBullet{\textbf{Prediction/control horizon:} adjust how far MPC plans ahead.}
  \SmallBullet{\textbf{Reference or residual term:} adding small correction action.}
  \SmallBullet{\textbf{Model multiplier:} adapt model online.}
\end{itemize}
}
\end{column}
\end{columns}
\end{frame}

%=========================================================
\begin{frame}
\PosterTitle{RL-Assisted MPC Results}
\scriptsize

\vspace{0.2cm}

\begin{columns}[T,onlytextwidth]
\begin{column}{0.65\textwidth}
\ColumnHeader{maccblue}{Polymer combined supervisor vs nominal MPC}
\vspace{0.05cm}
\FigurePanel{compare_rewards.png}{0.39\textheight}\\[-0.08em]

\vspace{0.05cm}
\FigurePanel{compare_outputs_last_episode.png}{0.39\textheight}\\[-0.08em]
\end{column}

\begin{column}{0.34\textwidth}
\ColumnHeader{maccgreen}{Reward gains vs nominal MPC}
\vspace{0.06cm}
\TextPanel{maccgreen}{softgreen}{0.3\textheight}{%
\centering
{\color{maccgreen}\bfseries Polymer}\\[-0.12em]
{\tiny
\setlength{\tabcolsep}{2.0pt}
\renewcommand{\arraystretch}{0.88}
\begin{tabular}{lcc}
\hline
Method & Full & Eval. ep. \\
\hline
Scalar matrix & +0.5175 & +0.8781 \\
Structured matrix & +0.4489 & +0.8870 \\
Residual TD3 & +0.4400 & +0.6912 \\
Weights TD3 & +1.5725 & +1.7743 \\
Horizon DQN & +1.5217 & +1.5827 \\
Dueling DQN & +1.5129 & +1.8686 \\
Combined & \textbf{+1.8891} & \textbf{+2.3761} \\
\hline
\end{tabular}}\\[-0.05em]
{\tiny Combined win rate: 98.0\%.}
}

\vspace{0.06cm}
\TextPanel{maccred}{softred}{0.15\textheight}{%
\centering
{\color{maccred}\bfseries C$_2$ splitter}\\[-0.12em]
{\tiny
\setlength{\tabcolsep}{2.0pt}
\renewcommand{\arraystretch}{0.90}
\begin{tabular}{lcc}
\hline
Method & Full & Eval. ep. \\
\hline
Dueling DQN & +2.1918 & +2.8352 \\
Weights SAC & +1.9945 & +0.3096 \\
\hline
\end{tabular}}\\[-0.05em]
{\tiny Only successful full-run distillation results shown.}
}

\vspace{0.06cm}
\ColumnHeader{maccblue}{Takeaway}
\vspace{0.05cm}
\TextPanel{maccblue}{softblue}{0.22\textheight}{%
\begin{itemize}
  \SmallBullet{The combined supervisor can select useful assistance when the operating regime changes.}
  \SmallBullet{Result: improved learning performance while keeping the MPC control structure.}
\end{itemize}
}
\end{column}
\end{columns}
\end{frame}

%=========================================================
\begin{frame}
\PosterTitle{Lyapunov-Filtered Safe RL}
\scriptsize

\vspace{0.2cm}

\begin{block}{\textbf{Safe deployment idea}}
\begin{itemize}\justifying
  \item TD3 can propose high-performance control moves, but a model-based safety layer decides what reaches the plant.
  \item If the RL move violates bounds or a Lyapunov contraction check, the supervisor falls back to MPC.
\end{itemize}
\end{block}

\vspace{-0.16cm}
\begin{columns}[T,onlytextwidth]
\begin{column}{0.50\textwidth}
\ColumnHeader{maccred}{Safe RL architecture}
\vspace{0.06cm}
\FigurePanel{SafeRL.png}{0.44\textheight}
\end{column}

\begin{column}{0.47\textwidth}
\ColumnHeader{maccred}{Safety stack}
\vspace{0.06cm}
\TextPanel{maccred}{softred}{0.4\textheight}{%
\begin{itemize}
  \SmallBullet{Target selector computes admissible local steady targets.}
  \SmallBullet{TD3 proposes a candidate action in deviation coordinates.}
  \SmallBullet{Lyapunov filter checks bounds, move limits, and contraction.}
  \SmallBullet{MPC fallback supplies the applied input when RL is rejected.}
\end{itemize}
}

\end{column}
\end{columns}

\vfill
\begin{center}
    \includegraphics[height=0.95cm]{MACC.png}\hspace{0.25cm}
    \includegraphics[height=0.95cm]{Linde.png}
\end{center}
\end{frame}

\end{document}
